\subsection{The Problem} 
Consider the polynomial
\[
x^5-3x^2-1
\]
It has two possible rational roots, \makeblank, and we can easily see that neither is a root:
\begin{lpic}{}{2}
\end{lpic}
But we can also use Descartes's Rule of Signs to analyze the roots:
\begin{itemize}
\item The polynomial has \makeblanksmall sign change, so it has \makeblanksmall positive root.
\item If we substitute $-x$ we get \makeblank which has \makeblanksmall sign changes, so the polynomial has \makeblanksmall negative roots.
\end{itemize}
So this polynomial has \makeblanksmall real root (and it must have \makeblank[1.5in] imaginary roots). If we have a computer generate a graph of the polynomial for us, we can see a zero:
\begin{icpic}
\setgraphsmall
\GraphLarge
\draw[graph,domain=-.8976:1.6828,samples=100] plot (\x,{(\x)^5-3*(\x)^2-1});
\end{icpic}
It looks like it's somewhere between \makeblanksmall and \makeblanksmall.
\begin{itemize}
\item The computer draws the graph by plotting lots of points and connecting them.
\item How do we know that there's \emph{really} a zero there?
\item How can we know what it is?
\end{itemize}